\section{Methods}
\label{sec:methods}

\subsection{Task}
The environment is a spatially cued orientation-change-detection task with a horizon of $T=29$ frames. Each $50\times50\times3$ frame is partitioned into a $10\times10$ grid of $5\times5$ patches ($N=100$ tokens). Stimuli are oriented patches whose orientations are corrupted by frame-to-frame Gaussian noise ($\sigma=5$) to set difficulty. A coloured cue is presented on a single early frame and then removed; its side indicates the likely change location and its colour sets the trial's reward magnitude. The cue's reliability (validity $\in\{0.25,0.50,0.75,1.00\}$) is the probability that an upcoming change, present on half of trials, occurs at the cued location. On change trials the target orientation shifts by a magnitude of up to roughly $64^{\circ}$, with onset jittered across frames (change index in roughly $[11,25]$) and the new orientation persisting to the end of the episode; the change magnitude is reduced over training by a performance-gated curriculum. On each frame the agent emits a binary wait/declare action: declaring at or after the change is a hit, declaring before it is a false alarm, and waiting out a no-change trial is a correct rejection. Reward is delivered for hits (scaled by the cue-colour value) and for correct rejections.

\subsection{Agent}
The agent is a conv-free recurrent vision transformer (Section~\ref{sec:architectures}). Each patch is embedded by a shared per-patch network into a $d=128$-dimensional token with learned positional and temporal codes; cross-location information flows only through attention. The encoder maintains two per-token recurrent memories, $\mathbf{H}_1$ (sensory) and $\mathbf{H}_2$ (deep), advanced by independent LSTM cells, and reads them through two single-head, pre-norm cross-attention streams---a priority stream ($\mu$) feeding the actor and a value stream ($q$) feeding the critic---differing only in key/value bank and residual (Equations~\eqref{eq:arch-mu}--\eqref{eq:arch-q}). The actor and the distributional critic are one-dimensional-convolutional heads over the token axis; the actor emits a wait/declare logit pair and the critic a five-quantile distribution over returns. The three configurations of the central comparison---V1 (shared memory), V2 (independent), V3 (swapped routing)---hold the front-end, the heads, the parameter budget and the entire training recipe fixed, differing only in inter-stream coupling and readout routing.

\subsection{Training}
Training uses an offline actor--critic objective with identical hyperparameters and identical update counts across the three configurations. The actor loss is a maximum-a-posteriori policy-improvement (MPO/MAP) term blended with a behaviour-cloning term; the critic is trained with a quantile Bellman loss against a periodically hard-copied target network that supplies both the bootstrap target and the policy-improvement reference. Whole episodes are stored in a prioritised replay buffer with priorities set by per-episode quantile error, and a brief forced-wait warm-up trains the critic on full trajectories before the actor begins to act.

\subsection{Evaluation}
Each trained configuration is evaluated on $500$ fresh episodes under a stochastic (sampled) policy and scored by signal-detection category (hit, miss, false alarm, correct rejection), from which we compute overall accuracy, hit rate on change trials, false-alarm rate on no-change trials, and reaction time on hits. For the working configuration, the cueing analysis splits hit rate by cue reliability and by cued versus uncued change location, with per-validity cells of $n=300$ and per-magnitude psychometric cells of $n=250$; reaction time is summarised by its median within condition. Across all $500$ episodes the two collapsed configurations declare a change exactly zero times, so their hit and false-alarm rates are both $0.000$ and no psychometric or reaction-time structure is defined. The three configurations of the central comparison correspond to internal model identifiers \texttt{v11\_part2} (V1, shared memory), \texttt{v11\_part4} (V2, independent) and \texttt{v11\_part3} (V3, swapped routing).
